\documentclass[a4paper,12pt]{ctexart}
\usepackage{xcolor}
\usepackage[top=1.8cm, bottom=1.8cm, left=1.8cm, right=1.8cm]{geometry}
\usepackage{array}
\linespread{1.35}

\title{\textcolor{blue!70!black}{\Huge\textbf{前后端开发入门}}}
\author{}
\date{}

\begin{document}
\maketitle
\thispagestyle{empty}

\section*{\textcolor{blue!70!black}{说明}}

本模块以完成一个小项目为主线，介绍网站前后端开发的基本概念。目标：学生能独立搭建一个简单的网页应用（如个人博客或留言板）。

\vspace{0.3cm}

\textbf{一、什么是前后端？}

\begin{itemize}
  \item \textbf{前端（Frontend）}：用户看到的页面，负责外观和交互。使用 HTML（结构）+ CSS（样式）+ JavaScript（行为）。
  \item \textbf{后端（Backend）}：服务器上的程序，负责处理数据和业务逻辑。常用语言包括 Python（Flask/Django）、JavaScript（Node.js）、Java、Go 等。
  \item \textbf{数据库（Database）}：存储用户数据的地方，比如用户名、密码、文章内容。
\end{itemize}

\vspace{0.3cm}

\textbf{二、小项目路线（初拟）：}

\begin{enumerate}
  \item \textbf{项目一：个人主页}——纯前端，用 HTML + CSS 写一个自我介绍页面
  \item \textbf{项目二：互动页面}——加上 JavaScript，实现按钮点击、表单验证
  \item \textbf{项目三：留言板}——引入后端（Python Flask）和数据库（SQLite），实现留言的提交和展示
  \item \textbf{项目四：简易博客}——综合练习，包含用户登录、文章发布、评论功能
\end{enumerate}

\vspace{0.3cm}

\textbf{三、前置要求：} 建议先完成 Python 基础模块，再学习前后端开发。

\end{document}
